\documentclass[11pt]{article}
\usepackage{fontspec}
\usepackage[french]{babel}
\usepackage{amsmath}
\usepackage{amssymb}
\usepackage{amsthm}
\usepackage{mathrsfs}
\usepackage{tikz-cd}
\usepackage{enumitem}
\usetikzlibrary{babel}

\newcommand{\printedpage}[1]{\marginpar{\footnotesize p.~#1}}
\newcommand{\msp}[1]{\marginpar{\footnotesize [#1]}}
\newcommand{\uFl}{\underline{\mathrm{Fl}}}
\newcommand{\Fl}{\mathrm{Fl}}
\newcommand{\Filt}{\mathrm{Filt}}
\newcommand{\cM}{\mathcal{M}}
\newcommand{\card}{\operatorname{card}}
\newcommand{\circled}[1]{\tikz[baseline=(c.base)]\node[draw,circle,inner sep=0.5pt](c){#1};}
% footnote mark printed as superscript "(n)"
\newcommand{\pfn}[2]{\textsuperscript{(}\footnote[#1]{#2}\textsuperscript{)}}

\begin{document}

\printedpage{9}
% head: A. Grothendieck : Dérivateurs, Ch. XV, Édition des Universités de Montpellier II et Paris VII
% foot: Version 5 novembre 2010

{\itshape d'où par itération un système inductif ordinal d'endofoncteurs $f\longmapsto p_\alpha(f)$ de $\uFl(\cM)$, et d'homomorphismes $i_\alpha : \mathrm{id}_{\uFl(\cM)} \longrightarrow p_\alpha$, \textup{i.e.}}
\[
\begin{tikzcd}[column sep=small]
X \arrow[rr, "i_\alpha(f)"] \arrow[dr, "f"'] & & X_\alpha(f) \arrow[dl, "p_\alpha(f)"] \\
& Y\,, &
\end{tikzcd}
\]
{\itshape formant un système inductif ordinal de factorisations de $f$ fonctorielles en $f$.}

\msp{page 11}

{\itshape Ceci posé, on a :}
\begin{enumerate}[label={}, leftmargin=*]
\item[a)] {\itshape Pour tout ordinal $\alpha$, on a $i_\alpha(f)\in\Phi$.}
\item[b)] {\itshape Supposons que les sources des $i\in C_0$ soient accessibles et soit alors $\pi$ un cardinal dans $\mathfrak{U}$ tel qu'elles soient $\pi$-accessibles. Alors pour tout nombre ordinal $\alpha$ grand devant $\pi$, on a $p_\alpha(f)\in C_{0*}$, donc}
\[
f \;=\; p_\alpha(f)\, i_\alpha(f) \qquad\quad \textit{avec} \qquad\quad i_\alpha(f)\in\Phi\,, \qquad\quad p_\alpha(f)\in\Psi\,.
\]
\end{enumerate}

On a un résultat légèrement plus précis, que voici :

\textbf{Corollaire 1.} \textit{Soit $\Phi_0$ une partie de} $\Fl(\cM)$ \textit{telle que}
\[
C_0''' \;\subset\; \Phi_0 \;\subset\; \Phi \qquad\quad (\textit{d'où}\ \ C_0\subset\Phi_0\subset\Phi)\,,
\tag*{(1.24)}
\]
\textit{où $C_0'''\subset\Phi$ est définie dans} 1.19, \textit{p.} 8. \textit{Par exemple, si $\Phi_0$ est une partie de} $\Fl(\cM)$ \textit{satisfaisant}
\[
C_0 \;\subset\; \Phi_0 \;\subset\; \Phi\,,
\]
\textit{et stable par cochangement de base, par petites sommes amalgamées} (\textit{cf.} (1.18), p.~8), \textit{par $\varinjlim$ ordinales strictes} (\textit{cf.} (1.19), p.~8), \textit{elle satisfait} (1.24). (\textbf{N.B.} \textit{Si $\cM$ admet des sommes amalgamées simples, de tels $\Phi_0$ existent trivialement, il suffit de prendre $\Phi_0=\Phi$.})

\msp{page 12}

\textit{Ceci posé, pour toute flèche $f$ dans $\cM$, et tout ordinal $\alpha$, on a}
\[
i_\alpha(f) \;\in\; \Phi_0\,.
\]

\textbf{Corollaire 2.} \textit{Supposons que les sources des flèches $i\in C_0$ soient accessibles, et soit $\Phi_0$ comme dans le corollaire} 1 (p.~ex. $\Phi_0=C_0'''$). \textit{Alors le couple $(\Phi_0,\Psi)$ est un précouple de Quillen,} i.e.
\begin{enumerate}[label={}, leftmargin=*]
\item[a)] $\Phi_0\longleftrightarrow\Psi$, i.e. $\Psi\subset\Phi_{0*}$ (\textit{ou ce qui revient au même, $\Phi_0\subset\Psi^*$}), \textit{et}
\item[b)] \textit{on a la propriété de factorisation de Quillen.}
\end{enumerate}
\textit{Cela implique donc que $\Phi$ $(=\widetilde{C_0}=\widetilde{\Phi_0})$ est formé des flèches qui sont facteurs directs de flèches de $\Phi_0$. Donc si $\Phi_0$ est stable par facteurs directs, on a $\Phi_0=\Phi$,} i.e. \textit{$\Phi_0$ est clos à gauche.}

\printedpage{10}
% head: A. Grothendieck : Dérivateurs, Ch. XV, Édition des Universités de Montpellier II et Paris VII
% foot: Version 5 novembre 2010

\section*{2\quad Variation des factorisations $f=p_\alpha i_\alpha$ avec la partie $C_0$ génératrice}

Je dis que le système inductif ordinal de foncteurs $p_\alpha=p_\alpha^{C_0} : \uFl(\cM)\longrightarrow\uFl(\cM)$,
\msp{page 13}
muni des flèches
\[
i_\alpha^{C_0} \;:\; \mathrm{id}_{\uFl(\cM)} \longrightarrow p_\alpha^{C_0}\,,
\]
dépend fonctoriellement de $C_0$, quand on considère $C_0$ comme un élément de l'ensemble ordonné des \textit{petites} parties de $\Fl(\cM)$, formées de flèches coquarrables (donc de l'ensemble ordonné
\[
\mathfrak{P}_{\mathfrak{U}\text{-petites}}(\Fl_{\text{coquar.}}(\cM)) \quad ).
\]
Si on a une inclusion
\[
C_0 \;\hookrightarrow\; C_1\,,
\tag*{(2.1)}
\]
on en déduit pour tout ordinal $\alpha$ un homomorphisme fonctoriel
\[
u_\alpha \;=\; u_\alpha^{C_1,C_0} \;:\; X_\alpha^{C_0} \;\longrightarrow\; X_\alpha^{C_1}\,,
\tag*{(2.2)}
\]
rendant commutatifs les carrés déduits de $f : X\longrightarrow Y$ dans $\cM$,
\[
\begin{tikzcd}
X \arrow[r, "i_\alpha^{C_0}(f)"] \arrow[d, "\mathrm{id}"'] & X_\alpha^{C_0}(f) \arrow[r, "p_\alpha^{C_0}(f)"] \arrow[d, "u_\alpha(f)"] & Y \arrow[d, "\mathrm{id}"] \\
X \arrow[r, "i_\alpha^{C_1}(f)"] & X_\alpha^{C_1}(f) \arrow[r, "p_\alpha^{C_1}(f)"] & Y\,,
\end{tikzcd}
\tag*{(2.3)}
\]
et compatible avec les morphismes de transition pour $\alpha$ variable,
\[
\begin{tikzcd}
X_\alpha^{C_0} \arrow[r, "{\mathrm{trans}_{\beta,\alpha}^{C_0}}"] \arrow[d, "u_\alpha"'] & X_\beta^{C_0} \arrow[d, "u_\beta"] \\
X_\alpha^{C_1} \arrow[r, "{\mathrm{trans}_{\beta,\alpha}^{C_1}}"] & X_\beta^{C_1}\,.
\end{tikzcd}
\tag*{(2.4)}
\]

\msp{page 14}

Pour définir les $u_\alpha$, satisfaisant aux conditions précédentes, il suffit de les définir pour $\alpha=1$ ``et d'itérer'' :
\[
\begin{tikzcd}[column sep=small]
& X_1^{C_0}(f) \arrow[dd, "u_1(f)"] \arrow[dr, "p_1^{C_0}(f)"] & \\
X \arrow[ur, "i_1^{C_0}(f)"] \arrow[dr, "i_1^{C_1}(f)"'] & & Y \\
& X_1^{C_1}(f) \arrow[ur, "p_1^{C_1}(f)"'] & {}.
\end{tikzcd}
\tag*{(2.5)}
\]
Or $X_1^{C_0}(f)$ est défini comme une somme directe dans $X\backslash\cM/Y$, relative à un ensemble d'indices $\Delta(C_0,f)$ \textit{contenu} dans $\Delta(C_1,f)$ -- et la famille d'objets dont on prend la somme directe étant induite d'une famille (non petite) dont l'ensemble d'indices serait $\Delta(\Fl_{\text{coquar.}}(\cM),f)$, savoir formé de \textit{tous} les diagrammes carrés
\printedpage{11}
% head: A. Grothendieck : Dérivateurs, Ch. XV, Édition des Universités de Montpellier II et Paris VII
% foot: Version 5 novembre 2010
\[
D : \qquad
\begin{tikzcd}
A \arrow[r, "u"] \arrow[d, "i"'] & X \arrow[d, "f"] \\
B \arrow[r, "v"] & Y
\end{tikzcd}
\qquad\qquad \text{avec $i$ coquarrable sans plus ,}
\]
la définition des $X_D$, $i_D$, $p_D$ associés à $D$ ne dépendant que de $D$, et non de l'ensemble $C_0$ ou $C_1$ de flèches coquarrables qui contient $i$. Ceci dit, la flèche $u_1(f)$
\msp{page 15}
dans (2.5) n'est autre que l'inclusion canonique d'une somme amalgamée partielle, indexée par $\Delta(C_0,f)$, dans la somme amalgamée correspondante à l'ensemble d'indices plus grand. La fonctorialité de cette flèche $u_1(f)$, pour $f$ variable, ``est immédiate'', elle résulte de la commutativité du carré, associé à un morphisme $v=(v_X,v_Y)$ de flèches
\[
\begin{tikzcd}
X \arrow[r, "f"] \arrow[d, "v_X"'] & Y \arrow[d, "v_Y"] \\
X' \arrow[r, "f'"] & Y'\,,
\end{tikzcd}
\tag*{(2.6)}
\]
savoir le diagramme d'applications entre ensembles d'indices
\[
\begin{tikzcd}
\Delta(C_0,f) \arrow[r, hook, "u_1(f)"] \arrow[d, "{\begin{array}{r}\text{dépendance fonctorielle}\\ \text{en $f$ de $\Delta(C_0,f)$}\end{array}}"'] & \Delta(C_1,f) \arrow[d, "{\begin{array}{l}\text{dépendance fonctorielle}\\ \text{en $f$ de $\Delta(C_1,f)$}\end{array}}"] \\
\Delta(C_0,f') \arrow[r, hook, "u_1(f')"] & \Delta(C_1,f')
\end{tikzcd}
\tag*{(2.7)}
\]
(commutativité qui revient à dire que $\Delta(C_0,f)$ peut être considéré comme un foncteur en le couple $(C_0,f)$ comme objet de la catégorie produit $\underline{\mathfrak{P}_{\text{pet.}}(\Fl_{\text{coquar.}}(\cM))}\times\uFl(\cM)$ \ldots). Le carré commutatif (2.7) est la partie ``indices'' d'un carré commutatif entre familles d'objets $\mathscr{X}(C_0,f)$, $\mathscr{X}(C_0,f')$, $\mathscr{X}(C_1,f)$, $\mathscr{X}(C_1,f')$,
\msp{page 16}
indexés par ces ensembles, savoir les $X_D^{C_0}$ et les $X_D^{C_1}$ relatifs aux diagrammes $D$ impliquant, soit la flèche $f$, soit la flèche $f'$ \ldots

Or la construction canonique du système inductif ordinal $p_\alpha$, avec les $i_\alpha : \mathrm{id}_{\uFl(\cM)}\longrightarrow p_\alpha$, est fonctorielle par rapport à la donnée initiale $p_1$, $i_1$ -- je n'explicite pas le détail de la fonctorialité, bien que je devrais \ldots

\section*{3\quad Majoration des cardinaux}

Je suppose à présent que $\cM$ est muni d'une \textit{filtration cardinale}
\[
(\Filt^\pi)_{\pi\geq\pi_0} \qquad\quad (\pi_0,\ \pi \text{ cardinaux dans } \mathfrak{U})
\tag*{(3.1)}
\]
(SGA 4 I 9.12, p.~149), les $\Filt^\pi$ stables par sommes amalgamées simples (chaque fois qu'elles existent), donc aussi par sommes amalgamées infinies indexées par des ensembles
\printedpage{12}
% head: A. Grothendieck : Dérivateurs, Ch. XV, Édition des Universités de Montpellier II et Paris VII
% foot: Version 5 novembre 2010
d'indices de cardinal $\leq\pi$. On suppose de plus que tout objet de $\Filt^{\pi_0}$ soit accessible, et on désigne par \circled{$\pi_1$} le plus petit cardinal tel que $\pi_1\geq\pi_0$ et que les objets de $\Filt^{\pi_0}$ soient $\pi_1$-accessibles. Désignant, pour tout cardinal, par $\cM_\pi$ le sous-ensemble de $\mathrm{Ob}\,\cM$ formé des objets $\pi$-accessibles, on aura alors pour tout $\pi\geq\pi_1$ (\textit{loc. cit.} 9.16)
\[
\cM_\pi \;\subset\; \Filt^\pi \;\subset\; \cM_{(\pi^{\pi_0})} \qquad\quad (\text{si } \pi\geq\pi_1)\,,
\tag*{(3.2)}
\]

\msp{page 17}

donc si $\pi^{\pi_0}=\pi$, p.~ex. si $\pi=2^c$ avec $c\geq\pi_0$, on aura
\[
\Filt^\pi \;=\; \cM_\pi\,.
\tag*{(3.2 bis)}
\]
(Si on admet l'hypothèse du continu généralisée, on a $\pi^{\pi_0}=\pi$ pour tout $\pi>\pi_0$, donc $\Filt^\pi=\cM_\pi$ si $\pi\geq\pi_1$ et $\pi>\pi_0$.)

Je pose, pour tout cardinal $\pi\geq\pi_0$,
\[
\Fl^\pi(\cM) \;=\; \{\ f\in\Fl(\cM) \mid \text{source et but de $f$ sont dans } \Filt^\pi\ \}
\tag*{(3.3)}
\]
-- ce qui sans doute donne une filtration cardinale sur $\uFl(\cM)$ (mais ce n'est pas là notre souci \ldots). Mon but à présent, c'est de dégager des conditions sur $C_0\subset\Fl_{\text{coquar.}}(\cM)$ (petite partie de $\Fl(\cM)$) et sur le cardinal $\pi$, et l'ordinal $\alpha$, pour que l'on ait
\[
\begin{aligned}
f\in\Fl^\pi(\cM) \;\;\overset{?}{\Longrightarrow}\;\; & X_\alpha(f)\in\Filt^\pi \\
& (\textit{i.e.}\ p_\alpha(f)\in\Fl^\pi(\cM),\ \text{ou encore } i_\alpha(f)\in\Fl^\pi(\cM))\,.
\end{aligned}
\tag*{(3.4)}
\]
Plus généralement, je me propose de donner une estimation de la ``cardinalité'' de $X_\alpha(f)$, pour $C_0$ fixé, en fonction de $C_0$, de $f$, de $\alpha$, plus exactement, en fonction des trois cardinaux $\pi(C_0)$, $\pi(f)$, $\pi(\alpha)$ définis comme
\[
\left\{
\begin{array}{rcl}
\pi(C_0) & = & \text{plus petit cardinal tel que } C_0\subset\Fl^\pi(\cM) \\[1ex]
\pi(f) & = & \pi(\{f\}) \;=\; \text{plus petit cardinal tel que } f\in\Fl^\pi(\cM) \\[1ex]
\pi(\alpha) & = & \card\alpha\,.
\end{array}
\right.
\tag*{(3.5)}
\]
Donc je cherche une fonction
\msp{page 18}
cardinale de variables cardinales $\varphi$, telle que l'on ait
\[
X_\alpha(f) \;\in\; \Filt^\pi \qquad\quad \text{si} \qquad\quad \pi \;=\; \varphi(\underbrace{\pi(C_0)}_{r},\underbrace{\pi(f)}_{s},\underbrace{\pi(\alpha)}_{t})\,.
\tag*{(3.6)}
\]
L'estimation cruciale concernera sûrement le cas $\alpha=1$. Notons que $X_1^{C_0}(f)=X_1(f)$ est défini comme une somme amalgamée (en général infinie) d'objets $X_D$, pour un ensemble $\Delta(C_0,f)=\Delta(f)$ de diagrammes, savoir les diagrammes commutatifs
\[
\begin{tikzcd}
A \arrow[r] \arrow[d, "i"'] & X \arrow[d, "f"] \\
B \arrow[r] & Y
\end{tikzcd}
\]
\printedpage{13}
% head: A. Grothendieck : Dérivateurs, Ch. XV, Édition des Universités de Montpellier II et Paris VII
% foot: Version 5 novembre 2010
avec $i\in C_0$. Pour un $i$ fixé, et un $f$ fixé, le cardinal de l'ensemble des diagrammes de ce type est majoré par celui de $\mathrm{Hom}(A,X)\times\mathrm{Hom}(B,Y)$. On a $A,B\in\Filt^r$, $X,Y\in\Filt^s$, donc il nous faut une majoration de
\[
\card\mathrm{Hom}_{\cM}(A,X)\,, \qquad\quad A\;\in\;\Filt^r\,, \qquad\quad X\;\in\;\Filt^s\,.
\]

\msp{page 19}

Voici les majorations essentielles dans ce sens :

\textbf{Proposition 3.1.} \textit{Soit $\pi_0$ le cardinal initial de la filtration cardinale} $(\Filt^\pi)_{\pi\geq\pi_0}$ \textit{envisagé, $\pi_1$ \underline{le plus petit} cardinal $\geq\pi_0$ tel que les objets de} $\Filt^{\pi_0}$ \textit{soient $\pi_1$-accessibles, $C_1=$} card $\underbrace{\widetilde{\uFl^{\pi_1}(\cM)}}_{=\uFl(\cM^{\pi_1})}$\pfn{4}{\textit{cf.} plus bas la notation $\widetilde{\mathcal{C}}$.}. \textit{Pour tout cardinal $\pi\geq\pi_0$, je désigne par $\cM^\pi$ la sous-catégorie strictement \underline{pleine} de $\cM$ dont l'ensemble des objets est} $\Filt^\pi$. \textit{Enfin, si $\mathcal{C}$ est une catégorie, je désigne par $\widetilde{\mathcal{C}}$ l'ensemble des classes d'isomorphie d'objets de $\mathcal{C}$.}\pfn{5}{Signalons qu'on a toujours $\card\widetilde{\uFl(\mathcal{C})}\geq\card\widetilde{\mathcal{C}}$ car $\widetilde{\mathcal{C}}\hookrightarrow\widetilde{\uFl(\mathcal{C})}$ via $\mathrm{cl}(X)\longmapsto\mathrm{cl}(\mathrm{id}_X)$.}
% ?? uncertain: straight underlines under "le plus petit" and under "pleine" -- underlined / not underlined (a wavy tilde and an underbrace stand close by)

$1^\circ$) \textit{Si $r\geq\sup(\pi_1,C_1)$, on a}
\[
\card\,(\underbrace{\widetilde{\uFl^r(\cM)}}_{=\uFl(\cM^r)}) \;\leq\; 2^{(r^{\pi_0})} \quad\ (=2^r,\ \textit{si } r=2^c,\ \textit{avec } c\geq\pi_0)\,.
\tag*{(3.7)}
\]

$2^\circ$) \textit{Si $r,s\geq\sup(\pi_1,C_1)$, et si $A\in\Filt^r$, $X\in\Filt^s$, on a}
% (3.8) and (3.9) share one brace on the right
\[
\begin{array}{l}
\text{(3.8)}\\[3ex]
\text{(3.9)}
\end{array}
\qquad\qquad
\begin{array}{r@{\;\;\leq\;\;}l}
\card\mathrm{Hom}_{\cM}(A,X) & s^{(r^{\pi_0})}\\[3ex]
\card\widetilde{\uFl(\cM^r/X)} & s^{(r^{\pi_0})}
\end{array}
\qquad\quad
\left\{
\begin{array}{lll}
= & s^r & \textit{si } r=2^c,\ \textit{avec } c\geq\pi_0\,,\\[3ex]
= & s & \textit{si } s=2^{c'},\ \textit{avec } c'\geq r^{\pi_0}\,.
\end{array}
\right.
\]

\textbf{N.B.} Si on admet l'hypothèse du continu généralisée, qui implique que $\pi^{\pi'}=\pi$ si $\pi>\pi'$ ($\pi$ cardinal infini, $\pi'\geq1$), alors prenant $r$ tel que $r>\pi_0$, les formules précédentes donnent (avec de plus $r\geq\sup(\pi_1,C_1)$) :
\[
\card\widetilde{\uFl^r(\cM)} \;\;\leq\;\; 2^r\,,
\tag*{(3.7 bis)}
\]
\[
\card\mathrm{Hom}_{\cM}(A,X) \;\leq\; s^r \qquad\quad (\text{si } A\in\Filt^r,\ X\in\Filt^s)\,,
\tag*{(3.8 bis)}
\]
\[
\card\widetilde{\uFl(\cM^r/X)} \;\leq\; s^r \qquad\quad (\text{si } X\in\Filt^s)\,.
\tag*{(3.9 bis)}
\]
Sans hypothèse du continu, ces formules sont aussi valables pour $r$ de la forme $2^c$, avec $c\geq\pi_0$. Ce sont ces dernières formules qui on l'air les plus naturelles. Pour (3.8 bis) et % [sic] qui on l'air -> qui ont l'air, verb form
(3.9 bis), je signale que le deuxième membre $s^r$ est le cardinal de l'ensemble $\mathrm{Hom}_{\mathrm{Ens}}(A,X)$, où $A$, $X$ sont deux ensembles tels que $\card A=r$, $\card X=s$. On ne peut donc espérer améliorer (3.8 bis)\,! Pour la

\end{document}
